% Extraído de Textuais/04-Resultados.tex
\definecolor{meuVerdeHex}{HTML}{159d15} 


\begin{figure}[htbp]
    \centering
    \caption{Correlação de Spearman entre a ordem prevista e a real das \textit{features}, por \textit{release} de teste}
    \label{fig:res-spearman}
    \begin{tikzpicture}
        \begin{axis}[width=\linewidth, height=6.5cm, xlabel={\textit{Release} de teste (índice global)},
            ylabel={Spearman}, xmin=36.5, xmax=53.5, ymin=0, ymax=1, grid=major, grid style={gray!25},
            legend style={font=\footnotesize, at={(0.5,-0.25)}, anchor=north, legend columns=3, draw=none},
            tick label style={font=\footnotesize}, label style={font=\footnotesize},
            /pgf/number format/use comma]
            \addplot[meuVerdeHex!70!black, thick, mark=*, mark size=1.2pt] table[x=release, y=modelo]
                {Textuais/dados/spearman_por_release.dat};
            \addlegendentry{SVM linear (model3, ROS)}
            \addplot[blue!70!black, thick, dashed, mark=square*, mark size=1.2pt] table[x=release, y=persistencia]
                {Textuais/dados/spearman_por_release.dat};
            \addlegendentry{Persistência}
            \addplot[gray, mark=triangle*, mark size=1.2pt] table[x=release, y=melhor_metodo]
                {Textuais/dados/spearman_por_release.dat};
            \addlegendentry{LightGBM (model2, ROS)}
        \end{axis}
    \end{tikzpicture}
    \fonte{Elaborado pelo autor (2026)}
\end{figure}%
